



































\newcommand{\smallstrut}{\vrule height 0.5\baselineskip depth 0pt width 0pt}

\begin{figure}[t]
\centering
\small
\begin{subfigure}[T]{0.295\textwidth}
\centering
\caption{\textbf{Dialogue:} $x$, $\mathbf{y_{t}}$}
\label{fig:dialog1}
\begin{minted}[fontsize=\footnotesize, frame=single,linenos=false,breaklines,breaksymbol=,escapeinside=||,bgcolor=Box1Color]{text}
User: I am interested in playing Table tennis.

Response: I'm sure it's a great way to socialize, stay active
\end{minted}
\end{subfigure}
\begin{subfigure}[T]{0.34\textwidth}
\centering
\caption{\textbf{\fbmod} $~~\mathbf{fb}$}
\label{fig:dialog2}
\begin{minted}[fontsize=\footnotesize, frame=single,linenos=false,breaklines,breaksymbol=, escapeinside=||,bgcolor=Box2Color]{text}
Engaging: Provides no information about table tennis or how to play it.

User understanding: Lacks understanding of user's needs and state of mind.
\end{minted}
\end{subfigure}
\begin{subfigure}[T]{0.351\textwidth}
\centering
\caption{\textbf{\itermod} $~~\mathbf{y_{t+1}}$}
\label{fig:dialog3}
\begin{minted}[fontsize=\footnotesize, frame=single,linenos=false,breaklines,breaksymbol=, escapeinside=||,bgcolor=Box3Color]{text}
|\textbf{Response (refined): That's great to hear (...) ! It's a fun sport requiring quick reflexes and good hand-eye coordination. Have you played before, or are you looking to learn?}|
\end{minted}
\end{subfigure}%
\vspace{0.2em}
\begin{subfigure}[T]{0.34\textwidth}
    \centering
    \caption{\textbf{Code optimization:} $x$, $\mathbf{y_{t}}$}
    \label{fig:code1}
    \setminted{bgcolor=Box1Color}
    \begin{minted}[fontsize=\footnotesize, frame=single,linenos=false,breaklines,breaksymbol=,escapeinside=||]{python}
|\texttt{Generate sum of 1, ..., N}|
def sum(n):
    res = 0
    for i in range(n+1):
        res += i
    return res
    \end{minted}
\end{subfigure}%
\hspace{0.025\textwidth}
\vspace{1em}
\begin{subfigure}[T]{0.28\textwidth}
\centering
\caption{\textbf{\fbmod} $~~\mathbf{fb}$}
\label{fig:code2}
\begin{minted}[fontsize=\footnotesize, frame=single,linenos=false,breaklines,breaksymbol=, escapeinside=||,bgcolor=Box2Color]{text}
|\textbf{This code is slow as it uses brute force. A better approach is to use the formula ... (n(n+1))/2.}
|
\end{minted}
\end{subfigure}
\hspace{0.025\textwidth}
\begin{subfigure}[T]{0.31\textwidth}
    \centering
    \caption{\textbf{\itermod} $~~\mathbf{y_{t+1}}$}
    \label{fig:code3}
    \setminted{bgcolor=Box1Color}
    \begin{minted}[fontsize=\footnotesize, breaklines, frame=single,bgcolor=Box3Color]{python}
Code (refined)

def sum_faster(n):
  return (n*(n+1))//2

  
    \end{minted}
\end{subfigure}%
\caption{Examples of \ours: an initial output \colorbox{Box1Color}{\smallstrut} generated by the base \llm and then passed back to the \textit{same} \llm to receive feedback \colorbox{Box2Color}{\smallstrut} to the \textit{same} \llm to refine the output \colorbox{Box3Color}{\smallstrut}. The top row illustrates this for dialog generation where an initial dialogue response can be transformed into a more engaging one that also understands the user by applying feedback. The bottom row illustrates this for code optimization where the code is made more efficient by applying feedback.
}
\label{fig:modules}
\end{figure}
